\emph{Summary:}
Denna modul bygger vidare på grundläggande filhantering med praktiska exempel
från verkligheten.
Studenter lär sig att validera filer, hantera fel, bearbeta textfiler med
dictionaries, och kombinera funktioner från flera moduler.
Modulen visar hur filhantering integreras med tidigare koncept som dictionaries,
loopar, och funktioner för att lösa verkliga problem.

\emph{Intended learning outcomes:}
Efter denna modul ska studenten kunna:
\begin{enumerate}
  \item \label{FilesMoreLO1}%
    Validera att filer existerar och hantera \mintinline{python}{FileNotFoundError}
    på lämpligt sätt.
  \item \label{FilesMoreLO2}%
    Bearbeta textfiler med dictionaries för räkning och analys
    (t.ex.\ ordfrekvens, bokstavsfrekvens).
  \item \label{FilesMoreLO3}%
    Komponera program genom att använda funktioner från flera moduler
    tillsammans.
  \item \label{FilesMoreLO4}%
    Transformera filinnehåll genom att läsa från en fil, bearbeta data,
    och skriva resultat till en annan fil eller samma fil.
  \item \label{FilesMoreLO5}%
    Skriva strukturerad output (CSV) från dictionary-baserad data för
    vidare analys.
\end{enumerate}

\emph{Prerequisites:}
Studenten bör ha genomfört grundläggande filhantering och känna till:
\begin{itemize}
  \item Grundläggande filoperationer (\mintinline{python}{open()},
    \mintinline{python}{read()}, \mintinline{python}{write()},
    \mintinline{python}{with}-satsen)
  \item Dictionaries och deras metoder
  \item Funktioner och moduler
  \item Loopar och if-satser
  \item Strängmanipulation
\end{itemize}

\emph{Reading:}
\dots
